\section{Future Work}
\label{sec:future-work}

The results presented in this thesis establish a  two-stage localisation architecture; several directions extend these capabilities toward full field deployment. The most immediate is outdoor validation of the precise-alignment stage under operational vineyard conditions. The SLAM pipeline was already evaluated on data collected in working vineyards; extending the registration pipeline to the same environment would characterise its behaviour under wind-induced vine motion, variable ambient lighting, and diverse cultivar geometries. Field trials incorporating an independent six-degree-of-freedom pose reference (e.g.\ a motion capture system or precision-surveyed fiducial markers) would establish absolute accuracy bounds, complementing the internal consistency metrics reported here.

The 100\% condition-level convergence observed for LiDAR motivates evaluation on a larger and more diverse set of vines spanning multiple vineyards, cultivars, and ages. Such testing would confirm the generality of the demonstrated low-centimetre accuracy and would provide the repeated measurements necessary to quantify run-to-run repeatability across the vine population.

A further direction is evaluating the LiDAR pipeline in vineyard conditions using the full 360° point cloud rather than the narrow angular aperture used in the indoor trials. Given that the LiDAR modality already achieves 100\% convergence with only a single-vine field of view, registering against a spatially extended 3DGS model spanning several bays and both adjacent rows would provide substantially more geometric constraint. This additional structure is expected to improve both convergence reliability and pose precision by exploiting the extended spatial baseline available in operational deployments.

On the computational side, the finding that all 35 hypothesis initialisations converge to the same solution for LiDAR registration enables a reduction to real-time operation: reducing the hypothesis count to 7--10 or parallelising the embarrassingly independent GICP runs across multiple cores would reduce pipeline latency from approximately 46\,s to 10--15\,s without sacrificing the demonstrated reliability.

The stereo pipeline uses FoundationStereo~\cite{wenFoundationStereoZeroShot2025} in a zero-shot configuration without domain-specific fine-tuning. On the KITTI~2015 leaderboard, fine-tuned stereo methods achieve approximately 2$\times$ lower error rates (e.g.\ D1-all of 1.28\% for fine-tuned methods versus 2.8\% for FoundationStereo zero-shot). Domain-specific fine-tuning of FoundationStereo on vineyard imagery therefore represents a direct path to reducing the stereo depth error component of registration RMSE. The existing 3DGS vineyard model could itself serve as a source of pseudo-ground-truth training data by rendering synthetic stereo pairs with known disparity, enabling a self-supervised adaptation loop without manual annotation.

The reference model itself offers room for improvement. The 3DGS models used in these experiments were reconstructed from single-sided imagery (front face of the vine only), leaving the rear surface unrepresented. Two-sided capture, which the research group has since adopted, provides correspondences from all sensor return directions, increases the inlier count available to ICP, and would likely improve registration accuracy and robustness.

End-to-end system integration testing by coupling the SLAM-derived coarse pose with precise-alignment refinement and downstream manipulator path planning would validate that the centimetre-scale registration accuracy demonstrated here translates to successful physical pruning operations under field conditions.